\documentclass[11pt]{article}
\usepackage[letterpaper,margin=1in]{geometry}
\usepackage[T1]{fontenc}
\usepackage[utf8]{inputenc}
\usepackage{lmodern,microtype}
\usepackage{amsmath,amssymb,mathtools,bm,slashed}
\usepackage{xcolor,enumitem,booktabs,array,needspace}
\usepackage[most]{tcolorbox}
\usepackage{fancyhdr,titlesec,parskip,hyperref}
\definecolor{qftblue}{HTML}{245A8D}
\definecolor{qftgreen}{HTML}{2F6B4F}
\definecolor{qftgray}{HTML}{F4F6F8}
\hypersetup{colorlinks=true,linkcolor=qftblue,urlcolor=qftblue,pdfauthor={Nima Lashkari}}
\pagestyle{fancy}\fancyhf{}
\lhead{PHYS 662: Quantum Field Theory I}\cfoot{\thepage}
\setlength{\headheight}{14pt}
\titleformat{\section}{\Large\bfseries\color{qftblue}}{\thesection.}{0.6em}{}
\titleformat{\subsection}{\large\bfseries\color{qftgreen}}{\thesubsection}{0.6em}{}
\setlist[enumerate]{leftmargin=2em,itemsep=0.25em,topsep=0.35em}
\setlist[itemize]{leftmargin=2em,itemsep=0.2em,topsep=0.35em}
\newcommand{\dd}{\mathrm d}
\newcommand{\ii}{\mathrm i}
\newcommand{\eps}{\varepsilon}
\newcommand{\tr}{\operatorname{tr}}
\newcommand{\Tr}{\operatorname{Tr}}
\newcommand{\diag}{\operatorname{diag}}
\newcommand{\Lag}{\mathcal L}
\newcommand{\D}{\mathcal D}
\newcommand{\vac}{\lvert0\rangle}
\newcommand{\vbra}{\langle0\rvert}
\newcommand{\Tord}{\mathrm T}
\newcommand{\ket}[1]{\lvert#1\rangle}
\newcommand{\bra}[1]{\langle#1\rvert}
\newtcolorbox{keyidea}[1][]{enhanced,breakable,colback=qftgray,colframe=qftblue,
 boxrule=0.7pt,arc=1.5mm,title={#1},fonttitle=\bfseries,coltitle=white}
\newtcolorbox{remarkbox}[1][]{enhanced,breakable,colback=white,colframe=qftgreen,
 boxrule=0.6pt,arc=1.5mm,title={#1},fonttitle=\bfseries,coltitle=white}
\newcommand{\module}[1]{\clearpage\section{#1}}
\newcommand{\lecturefront}[4]{%
\rhead{Lecture #1}\hypersetup{pdftitle={PHYS 662 Lecture #1: #2}}
\begin{document}\pagenumbering{gobble}
\begin{titlepage}\centering\vspace*{1.05in}
{\scshape\Large PHYS 662 --- Quantum Field Theory I\par}
\vspace{0.35in}{\Huge\bfseries\color{qftblue}Lecture #1\par}
\vspace{0.2in}{\LARGE\bfseries #2\par}
\vspace{0.2in}{\large\color{qftgreen}#3\par}
\vspace{0.6in}{\large Nima Lashkari\par}{\large Fall 2026\par}
\vfill\begin{minipage}{0.86\textwidth}\small #4\par\medskip
Course page: \url{https://nimalashkari.com/quantum-field-theory-2026-part-i/}
\end{minipage}\vspace{0.5in}\end{titlepage}
\pagenumbering{roman}\setcounter{tocdepth}{1}\tableofcontents
\vfill\begin{remarkbox}[Place in the course]
Lectures 17--18 develop the quantum-mechanical construction. Lecture 19
extends Gaussian integration and source methods to scalar fields.
Lecture 20 introduces Grassmann integration and fermionic fields.
Problem Set 6 brings these four steps together.
\end{remarkbox}\clearpage\pagenumbering{arabic}}
\lecturefront{20}{Path Integrals for Fermionic Fields}
{Coherent states, Grassmann sources, and the fermion-loop sign}
{Fermionic paths carry anticommuting values. We build their integration
rules from a single mode, derive the Dirac Gaussian with explicit source
conventions, and trace the signs in time ordering, thermal boundary
conditions, and a closed fermion loop.}

\section{Why coherent states are the useful starting point}
Use \(\hbar=c=1\) and the mostly-minus metric of Lectures 15--19.
A fermionic annihilation operator cannot have an arbitrary ordinary
complex eigenvalue: from \(c^2=0\), an eigenvalue would obey
\(\xi^2=0\). Grassmann numbers will provide such labels. First recall
how a coherent-state path integral works for a bosonic oscillator.

For \([a,a^\dagger]=1\), define unnormalized states
\begin{equation}
 |z\rangle=e^{za^\dagger}|0\rangle,\qquad
 \langle\bar z|=\langle0|e^{\bar za},\qquad
 a|z\rangle=z|z\rangle.
\end{equation}
Their overlap and resolution of the identity are
\begin{equation}
 \langle\bar z'|z\rangle=e^{\bar z'z},\qquad
 1=\int\frac{\dd^2z}{\pi}e^{-\bar zz}|z\rangle\langle\bar z|.
                                                        \label{eq:bosonidentity}
\end{equation}
Here \(\bar z=z^*\) on the integration contour. Expanding in number
states proves the identity, using
\(\int\dd^2z\,e^{-|z|^2}z^n\bar z^m/\pi=n!\delta_{nm}\).

If \(H=:h(a^\dagger,a):\) is normal ordered, then
\begin{equation}
 \langle\bar z_{j+1}|e^{-\ii\epsilon H}|z_j\rangle
 =e^{\bar z_{j+1}z_j}
 [1-\ii\epsilon h(\bar z_{j+1},z_j)+O(\epsilon^2)].
\end{equation}
The product of overlaps and Gaussian measures produces a first-order
kinetic term \(\ii\bar z\dot z\). The normal symbol uses labels from
adjacent time slices; replacing it by another symbol without correcting
the discretization can change the Hamiltonian.

The continuum expression for a matrix element between unnormalized
coherent states contains an endpoint factor:
\begin{equation}
 K(\bar z_f,z_i)=\int\D\bar z\D z\,
 \exp\!\left\{\bar z_fz(t_f)+\ii\int_{t_i}^{t_f}\dd t\,
 [\ii\bar z\dot z-h(\bar z,z)]\right\}.             \label{eq:bosonpath}
\end{equation}
Fix \(z(t_i)=z_i\) and \(\bar z(t_f)=\bar z_f\); the discrete
construction defines the remaining endpoint prescription. Treating
both variables as independently fixed at both ends would overconstrain
the first-order equations.

For an oscillator, \(z=\sqrt{m\omega/2}\,q+\ii p/\sqrt{2m\omega}\).
The Hamiltonian \(\omega(a^\dagger a+1/2)\) has symbol
\(\omega(\bar zz+1/2)\). The vacuum constant is not generated by
silently replacing this symbol with \(\omega\bar zz\).

\module{Grassmann algebra and Berezin integration}
Independent odd generators satisfy
\begin{equation}
 \theta_a\theta_b=-\theta_b\theta_a,\qquad \theta_a^2=0.
\end{equation}
Every function is a finite polynomial in finitely many generators.
Even elements commute with all elements, whereas two homogeneous odd
elements anticommute. Barred variables are independent generators;
the bar is a label, not a constraint imposed during integration.

Use \emph{left} derivatives throughout. Their product rule is
\begin{equation}
 \partial^L_\theta(FG)=(\partial^L_\theta F)G
 +(-1)^{|F|}F(\partial^L_\theta G),                   \label{eq:leibniz}
\end{equation}
for homogeneous parity \(|F|\). Thus
\(\partial^L_\theta(\bar\theta\theta)=-\bar\theta\),
whereas \(\partial^L_{\bar\theta}(\bar\theta\theta)=\theta\).
Odd derivatives anticommute. Write their order explicitly, remembering
that the rightmost derivative acts first.

Berezin integration is defined by
\begin{equation}
 \int\dd\theta\,1=0,\qquad \int\dd\theta\,\theta=1.
\end{equation}
Choose the paired orientation
\begin{equation}
 \int\dd\bar\theta\,\dd\theta\,\theta\bar\theta=1,
 \qquad
 \int\dd\bar\theta\,\dd\theta\,\bar\theta\theta=-1.
                                                        \label{eq:measure}
\end{equation}
Then
\begin{equation}
 \int\dd\bar\theta\dd\theta\,e^{-\bar\theta a\theta}
 =\int\dd\bar\theta\dd\theta(1-a\bar\theta\theta)=a.
                                                        \label{eq:onegauss}
\end{equation}
Unlike an ordinary Gaussian, this integral produces a positive power
of the quadratic coefficient. It is an algebraic operation and
requires no convergence condition on \(a\).

Translation leaves the measure invariant. Consequently the integral
of a total Grassmann derivative vanishes, the starting point for
fermionic Schwinger--Dyson identities. Under \(\theta'=a\theta\),
one instead has \(\dd\theta'=a^{-1}\dd\theta\): the Jacobian is
inverse to the ordinary one, as follows by preserving
\(\int\dd\theta'\theta'=1\).

\begin{keyidea}[Signs require a complete convention]
An integration order, a field order, and a left/right derivative
convention must be specified together. Changing a convention is
allowed, but changing just one displayed sign is not. Equations
\eqref{eq:leibniz}--\eqref{eq:measure} fix the conventions used below.
\end{keyidea}

\module{Fermionic coherent states and time slicing}
For \(\{c,c^\dagger\}=1\), let odd labels anticommute with \(c,c^\dagger\).
Choose
\begin{equation}
 |\xi\rangle=e^{-\xi c^\dagger}|0\rangle,\qquad
 \langle\bar\xi|=\langle0|e^{-c\bar\xi}.
\end{equation}
Expanding the finite exponentials gives
\begin{equation}
 c|\xi\rangle=\xi|\xi\rangle,\qquad
 \langle\bar\xi|c^\dagger=\langle\bar\xi|\bar\xi,
 \qquad \langle\bar\xi|\xi\rangle=e^{\bar\xi\xi}.
\end{equation}
With the measure of \eqref{eq:measure},
\begin{equation}
 1=\int\dd\bar\xi\dd\xi\,e^{-\bar\xi\xi}
 |\xi\rangle\langle\bar\xi|.                         \label{eq:fermidentity}
\end{equation}
For a direct check, expand the outer product in the states
\(|0\rangle,|1\rangle=c^\dagger|0\rangle\). The integrated coefficient
of each diagonal projector is one; the odd off-diagonal terms integrate
to zero. The same construction extends to several modes with a fixed
mode order.

For a normal-ordered even Hamiltonian, inserting \eqref{eq:fermidentity}
between \(N\) short evolution steps gives
\begin{align}
 K(\bar\xi_f,\xi_i)
 =\lim_{N\to\infty}\int\prod_{j=1}^{N-1}\dd\bar\xi_j\dd\xi_j\,
 \exp\Bigg[&-\sum_{j=1}^{N-1}\bar\xi_j\xi_j
 +\sum_{j=0}^{N-1}\bar\xi_{j+1}\xi_j\nonumber\\[-2pt]
 &-\ii\epsilon\sum_{j=0}^{N-1}h(\bar\xi_{j+1},\xi_j)\Bigg],
                                                        \label{eq:discreteferm}
\end{align}
where \(\xi_0=\xi_i\) and \(\bar\xi_N=\bar\xi_f\). The short-step
formula is accurate to the order needed for the limit. The overlap
terms combine into
\(\bar\xi_f\xi_{N-1}-\sum_{j=1}^{N-1}\bar\xi_j(\xi_j-\xi_{j-1})\).
Thus the continuum notation is
\begin{equation}
 K=\int\D\bar\xi\D\xi\,
 e^{\bar\xi_f\xi(t_f)+\ii\int\dd t\,[\ii\bar\xi\dot\xi-h(\bar\xi,\xi)]}.
                                                        \label{eq:fermpath}
\end{equation}
Only \(\xi(t_i)\) and \(\bar\xi(t_f)\) are prescribed. The endpoint
factor is essential for the first-order variational problem and for
the correct coherent-state matrix element.

For \(H=\omega c^\dagger c\), direct evaluation in the two-dimensional
Hilbert space gives
\begin{equation}
 K(\bar\xi_f,\xi_i;T)=1+\bar\xi_f\xi_i e^{-\ii\omega T}
 =e^{\bar\xi_f\xi_i e^{-\ii\omega T}}.
\end{equation}
It reduces to the overlap at \(T=0\). This is a useful exact check
of the kinetic sign and the boundary factor in \eqref{eq:fermpath}.

\module{The Grassmann Gaussian gives a determinant}
For \(n\) pairs, orient the measure so that the product of the
\(n\) monomials \(\theta_a\bar\theta_a\) integrates to one.
For an invertible ordinary matrix \(M\),
\begin{equation}
 \int\D\bar\theta\D\theta\,e^{-\bar\theta M\theta}=\det M.
                                                        \label{eq:det}
\end{equation}
Only the term containing every integration variable survives.
Its antisymmetrized coefficient is
\(\sum_\sigma\operatorname{sgn}(\sigma)\prod_a M_{a\sigma(a)}\),
which is the determinant. For two pairs this is
\(M_{11}M_{22}-M_{12}M_{21}\); the minus sign follows directly
from exchanging odd variables.

With independent odd sources \(\bar\eta,\eta\), completion of the square
works provided their order is retained:
\begin{align}
 &\int\D\bar\theta\D\theta\,
 e^{-\bar\theta M\theta+\bar\eta\theta+\bar\theta\eta}
   =\det M\,e^{\bar\eta M^{-1}\eta},                 \label{eq:euclgauss}\\
 &-\bar\theta M\theta+\bar\eta\theta+\bar\theta\eta
 =-(\bar\theta-\bar\eta M^{-1})M(\theta-M^{-1}\eta)
   +\bar\eta M^{-1}\eta.
\end{align}
Translation invariance of Berezin integration proves the first line.
No complex conjugacy condition relates the shifted barred and unbarred
variables: they are independent integration coordinates.

The normalized two-point integral is
\begin{equation}
 \frac{\int\D\bar\theta\D\theta\,
             \theta_a\bar\theta_b e^{-\bar\theta M\theta}}
      {\det M}=(M^{-1})_{ab}.
                                                        \label{eq:finitecov}
\end{equation}
For one pair, its numerator is one and its denominator is \(M\),
which verifies the orientation. Reversing the insertion order gives
minus the transposed-index expression.

\begin{keyidea}[Three Gaussian measures]
At a finite regulator, a real commuting Gaussian gives
\((\det M)^{-1/2}\), a complex commuting Gaussian gives
\((\det M)^{-1}\), and a complex Grassmann Gaussian gives
\(\det M\), up to source-independent measure constants.
Their covariances are inverses; their vacuum determinants differ.
\end{keyidea}
Fermionic nilpotence removes a convergence question at fixed dimension.
It does not remove ultraviolet divergences when the number of field
modes tends to infinity. A Dirac determinant still requires a regulator.

\module{The Dirac functional and ordered source derivatives}
Apply the coherent-state construction to every regulated spinor mode.
The canonical Hamiltonian has one-particle operator
\(-\ii\bm\alpha\cdot\bm\nabla+\beta m\), with
\(\alpha^j=\gamma^0\gamma^j\) and \(\beta=\gamma^0\).
Combining it with the coherent kinetic term gives
\(\ii\psi^\dagger\dot\psi-
\psi^\dagger(-\ii\bm\alpha\cdot\bm\nabla+\beta m)\psi
=\bar\psi(\ii\slashed\partial-m)\psi\).
Consequently the independent Grassmann fields have action
\begin{equation}
 S_0=\int\dd^4x\,\bar\psi D\psi,\qquad
 D=\ii\gamma^\mu\partial_\mu-m.
\end{equation}
In operator language \(\bar\psi=\psi^\dagger\gamma^0\); in the
functional integral its components are independent odd variables.
Vacuum boundary conditions select the Feynman inverse, so
\begin{equation}
 S_F=\ii R_D,\qquad
 S_F(x-y)=\int\frac{\dd^4p}{(2\pi)^4}
 \frac{\ii(\slashed p+m)e^{-\ii p\cdot(x-y)}}{p^2-m^2+\ii0}.
                                                        \label{eq:SF}
\end{equation}
Here \(R_D\) is the Feynman inverse of the Dirac operator,
with \(DR_D=1\). Keep it distinct from the scalar propagator
\(D_F\). The two physical propagators obey
\(S_F=(\ii\slashed\partial+m)D_F\) using the scalar \(D_F\) of Lecture 19.

Introduce odd sources in the fixed order
\begin{equation}
 Z_0[\bar\eta,\eta]=
 \frac{\int\D\bar\psi\D\psi\,
 e^{\ii(\bar\psi D\psi+\bar\eta\psi+\bar\psi\eta)}}
 {\int\D\bar\psi\D\psi\,e^{\ii\bar\psi D\psi}}
 =e^{-\ii\bar\eta R_D\eta}=e^{-\bar\eta S_F\eta}.     \label{eq:fermZ}
\end{equation}
The shifts are \(\psi\mapsto\psi+R_D\eta\) and
\(\bar\psi\mapsto\bar\psi+\bar\eta R_D\). The determinant
\(\det(-\ii D)\) cancels in this normalized free functional.

With all derivatives acting from the left,
\begin{equation}
 \frac{\delta^L Z}{\delta\bar\eta_a(x)}
 =\ii\langle\psi_a(x)\rangle_{\rm unnorm},\qquad
 \frac{\delta^L Z}{\delta\eta_b(y)}
 =-\ii\langle\bar\psi_b(y)\rangle_{\rm unnorm}.
\end{equation}
The bracket here includes the source exponential and the zero-source
normalization, not division by \(Z[\bar\eta,\eta]\). In particular,
\begin{equation}
 \langle\Tord\psi_a(x)\bar\psi_b(y)\rangle
 =\left.\frac{\delta^L}{\delta\bar\eta_a(x)}
           \frac{\delta^L}{\delta\eta_b(y)}Z_0\right|_0
 =S_{F,ab}(x-y).                                     \label{eq:derivatives}
\end{equation}
Reversing these two derivatives gives \(-S_{F,ab}\).
Indeed differentiating \(-\bar\eta S_F\eta\) first with respect to
\(\eta\) gives \(+\bar\eta S_F\). This one-line check fixes the sign.

\module{Fermionic Wick contractions and the source equation}
Fermionic time ordering includes the sign of the permutation:
\begin{equation}
 \Tord\psi_a(x)\bar\psi_b(y)=
 \theta(x^0-y^0)\psi_a(x)\bar\psi_b(y)
 -\theta(y^0-x^0)\bar\psi_b(y)\psi_a(x).
\end{equation}
Expanding the Gaussian source functional gives Wick's theorem with
the same permutation signs. Writing \(S_{ab}\) for the appropriate
spacetime and spinor kernel, two useful orderings are
\begin{align}
 \langle\Tord\psi_a\bar\psi_b\psi_c\bar\psi_d\rangle
 &=S_{ab}S_{cd}-S_{ad}S_{cb},                        \label{eq:wickinterleaved}\\
 \langle\Tord\psi_a\psi_b\bar\psi_c\bar\psi_d\rangle
 &=S_{ad}S_{bc}-S_{ac}S_{bd}.                         \label{eq:wickgrouped}
\end{align}
The relative minus sign is fixed by exchanging odd fields. The
overall form also depends on the stated external field order; one
cannot copy a bosonic pairing formula and append an arbitrary sign.

Let \(j^\mu(x)=\bar\psi(x)\gamma^\mu\psi(x)\). Removing the product
of one-point expectations, the connected contraction is
\begin{equation}
 \langle\Tord j^\mu(x)j^\nu(y)\rangle_c
 =-\tr[\gamma^\mu S_F(x-y)\gamma^\nu S_F(y-x)].       \label{eq:current}
\end{equation}
The trace closes the spinor indices around the loop. The minus sign
is the exchange sign left over from closing the fermion chain.
Coincident composite operators are understood with a regulator.

\subsection{Schwinger--Dyson check}
A total left derivative with respect to \(\bar\psi\) gives
\begin{equation}
 \left[\frac1\ii D_{ab}(x)
       \frac{\delta^L}{\delta\bar\eta_b(x)}+\eta_a(x)\right]Z_0=0.
                                                        \label{eq:fermSD}
\end{equation}
The source term order is as displayed. Since
\(\delta^L Z_0/\delta\bar\eta=-S_F\eta Z_0\), this identity is
equivalent to
\begin{equation}
 (\ii\slashed\partial_x-m)S_F(x-y)=\ii\delta^{(4)}(x-y)I_4.
\end{equation}
It agrees with the canonical result because differentiating fermionic
time ordering produces the equal-time anticommutator. In momentum
space, \((\slashed p-m)(\slashed p+m)=p^2-m^2\); in position
space, \((\ii\slashed\partial-m)(\ii\slashed\partial+m)
=-(\Box+m^2)\). This also verifies the scalar-to-Dirac relation.

\module{Thermal fermions and antiperiodicity}
A trace is not obtained by imposing the bosonic coherent-state closure
without a sign. For an even operator \(O\), one mode obeys
\begin{equation}
 \Tr O=\int\dd\bar\xi\dd\xi\,e^{-\bar\xi\xi}
 \langle-\bar\xi|O|\xi\rangle.                      \label{eq:trace}
\end{equation}
This identity follows by expanding in the empty and occupied states.
For \(O=1\), the integrand is
\((1-\bar\xi\xi)^2=1-2\bar\xi\xi\), whose integral is two.
Omitting the minus sign in the bra would give zero, the supertrace
\(\Tr(-1)^F\), rather than the ordinary trace.

For \(H=\omega c^\dagger c\), the coherent kernel of Section 3 gives
\begin{equation}
 Z_\beta=\int\dd\bar\xi\dd\xi\,e^{-\bar\xi\xi}
 e^{-\bar\xi\xi e^{-\beta\omega}}
 =1+e^{-\beta\omega}.
\end{equation}
Its two terms are the empty and occupied states. Inserting intermediate
states around the trace moves the sign in \eqref{eq:trace} to the seam
of the Euclidean time circle. Therefore
\begin{equation}
 \psi(\tau+\beta,\bm x)=-\psi(\tau,\bm x),\qquad
 \bar\psi(\tau+\beta,\bm x)=-\bar\psi(\tau,\bm x).
                                                        \label{eq:AP}
\end{equation}
The allowed frequencies are \(\omega_n=(2n+1)\pi/\beta\). Bosonic
thermal fields instead have \(2n\pi/\beta\). These boundary conditions
are imposed by the trace and statistics, not by a change in the local
classical field equation.

\subsection{The Euclidean Dirac operator}
Choose \(\gamma_E^4=\gamma^0\), \(\gamma_E^j=-\ii\gamma^j\).
They satisfy \(\{\gamma_E^a,\gamma_E^b\}=2\delta^{ab}\). Then
\begin{equation}
 S_E=\int\dd^4x_E\,\bar\psi D_E\psi,\qquad
 D_E=\gamma_E^a\partial_a+m,
\end{equation}
and the Euclidean functional is \(\det D_E\) at zero source, up to
normalization. On infinite Euclidean space,
\begin{equation}
 S_E^{(2)}(x)=\int\frac{\dd^4p}{(2\pi)^4}e^{\ii p\cdot x}
 \frac{-\ii\gamma_E^ap_a+m}{p^2+m^2}.
\end{equation}
At finite temperature replace the energy integral by the antiperiodic
frequency sum. The superscript distinguishes this two-point kernel
from the Euclidean action. In both cases barred fields remain
independent integration variables.

\module{Closed loops and a first look at a gauge interaction}
Treat \(A_\mu(x)\) as a prescribed external field and choose
\begin{equation}
 \mathcal L=\bar\psi(\ii\slashed\partial-m)\psi
 -e\bar\psi\gamma^\mu\psi A_\mu.
\end{equation}
This corresponds to \(D_\mu=\partial_\mu+\ii eA_\mu\); the vertex
from \(e^{\ii S}\) is \(-\ii e\gamma^\mu\). There is no need to
quantize the gauge field to see the fermion-loop sign.

Integrating the fermions gives the vacuum response
\begin{align}
 \frac{\mathcal Z[A]}{\mathcal Z[0]}
 &=\frac{\det(D-e\slashed A)}{\det D},\\
 \log\frac{\mathcal Z[A]}{\mathcal Z[0]}
 &=\Tr\log(1-eR_D\slashed A)
 =-\sum_{n\geq1}\frac{e^n}{n}\Tr[(R_D\slashed A)^n]. \label{eq:looplog}
\end{align}
\(\Tr\) includes spacetime and spinor indices. Each term is a cyclic
chain of propagators and insertions. A commuting complex Gaussian
has the inverse determinant and the opposite logarithm. This is the
origin of one extra minus sign for each closed fermion loop.

For illustration, define \(\ii\Pi^{\mu\nu}(q)\) as the amputated
two-background-field insertion. Its loop integrand is
\begin{equation}
 \ii\Pi^{\mu\nu}(q)=-( -\ii e)^2
 \int\frac{\dd^4k}{(2\pi)^4}
 \tr\!\left[
 \gamma^\mu\frac{\ii(\slashed k+m)}{k^2-m^2+\ii0}
 \gamma^\nu\frac{\ii(\slashed k+\slashed q+m)}{(k+q)^2-m^2+\ii0}
 \right].                                            \label{eq:polarization}
\end{equation}
The leading minus is the closed-loop sign; the other factors come
from the two vertices and two propagators. The \(1/2!\) in the
interaction expansion is canceled by the two assignments of labeled
external insertions. Evaluating this ultraviolet-divergent integral
requires a regulator and renormalization, beyond this lecture.

This object contributes to vacuum polarization. It is not a claim
that an isolated on-shell photon spontaneously decays into a real
electron--positron pair. Its external energy and momentum, analytic
structure, and eventual gauge-field dynamics must be treated separately.

\subsection*{Questions to check understanding}
\begin{enumerate}
\item Why do coherent-state amplitudes fix different variables at the two ends?
\item Verify the one-pair Gaussian and the derivative order in \eqref{eq:derivatives}.
\item Derive the relative sign in \eqref{eq:wickinterleaved} by one exchange.
\item How does the ordinary thermal trace enforce antiperiodicity?
\item Where do the determinant and closed-loop signs enter \eqref{eq:looplog}?
\end{enumerate}
\end{document}
